\documentclass[11pt,a4paper]{article}
\usepackage[utf8]{inputenc}
\usepackage[T1]{fontenc}
\usepackage[english]{babel}
\usepackage{amsmath,amssymb,booktabs}
\usepackage{geometry,hyperref}
\geometry{margin=2.7cm}
\title{Universal Minimum Curvature Hypothesis II:\\candidate classical dynamics and no-go tests}
\author{Pepe Donato\\Independent Researcher\\\texttt{donato.pepe.it@gmail.com}}
\date{Audited comparative draft --- 2026}
\begin{document}
\maketitle

\begin{abstract}
We compare three implementations of hypothetical scale $\kappa_0>0$: hard pointwise constraint, curvature barrier, and a separate coarse-grained RMS bound. Reproducible checks cover KKT algebra, elementary barrier properties, and RMS non-equivalence. Neither pointwise dynamics yet has complete variation and canonical structure. States: A and B \texttt{INCOMPLETE}; C \texttt{NON\_IDENTIFIABLE} and \texttt{ALTERNATIVE\_HYPOTHESIS}; core hypothesis \texttt{UNPROVEN}; no-go \texttt{NO\_GO\_NOT\_ESTABLISHED}.
\end{abstract}

\section{Scientific status}\label{sec:p2-status}
Paper I showed that an ideal timelike geodesic has $\kappa=0$. Paper II assumes no solution: it compares mechanisms and preserves negative findings. Geometric-action literature supplies variational and Hamiltonian methods \cite{ArreagaCapovillaGuven2001,CapovillaGuvenRojas2002}; acceleration-dependent models do not establish minimum curvature \cite{NesterenkoEtAl1995}.

\section{Method and levels}\label{sec:p2-method}
We distinguish \texttt{KINEMATIC}, \texttt{SYMBOLIC}, algebraic \texttt{CONSTRAINT}, derivative-count \texttt{VARIATIONAL}, and \texttt{CONJECTURAL} parts. Covariance, boundary terms, constraints, conservation, stability, causality, standard limit, and observables are separate gates.

\section{Candidate A: hard constraint}\label{sec:p2-a}
Consider
\begin{equation}\label{eq:p2-hard}
S_A=S_0+\int ds\,\lambda g,\qquad g=\kappa_0-\kappa\le0,
\end{equation}
with $\lambda\ge0$ and $\lambda g=0$. The preregistered check identifies interior, active-boundary, and infeasible branches. For every $\kappa_0>0$, data with $\kappa=0$ lie outside the feasible set. At $\kappa_0=0$, only set-level admissibility returns. Nonsmooth variation, active-set evolution, and Dirac classification are missing: state \texttt{INCOMPLETE}.

\section{Candidate B: barrier}\label{sec:p2-b}
We fix
\begin{equation}\label{eq:p2-barrier}
S_B=-mc\int ds+\epsilon mc\int ds\,\frac{1}{\kappa/\kappa_0-1},
\qquad \kappa>\kappa_0.
\end{equation}
Barrier diverges at boundary, decreases, and is convex. Limit is nonuniform: it vanishes at fixed $\kappa>0$, remains finite for $\kappa=r\kappa_0$, and excludes geodesic. Generic variation can contain up to four embedding derivatives, but degeneracy and gauge constraints prevent automatic instability conclusions. State \texttt{INCOMPLETE}.

\section{Candidate C: separate RMS}\label{sec:p2-c}
As a distinct proposal define
\begin{equation}\label{eq:p2-rms}
\kappa_{\rm RMS}=\left(\frac{\int d\tau\,w\kappa^2}{\int d\tau\,w}\right)^{1/2}\ge\kappa_0.
\end{equation}
For samples $[0,2]$, equal weights, and $\kappa_0=1$, RMS is $\sqrt2$ while pointwise bound fails. Thus C is \texttt{ALTERNATIVE\_HYPOTHESIS}, non-equivalent, and \texttt{NON\_IDENTIFIABLE} without kernel and instrument response.

\section{Dimensions}\label{sec:p2-dimensions}
$[\kappa]=[\kappa_0]=L^{-1}$ and $[S]=ML^2T^{-1}$. In B, $\epsilon$ and ratio are dimensionless. In A, dimension of $\lambda$ depends on normalized constraint $1-\kappa/\kappa_0$ versus dimensional $\kappa_0-\kappa$; convention must precede variation.

\section{Constraints, stability, and causality}\label{sec:p2-stability}
No complete canonical structure has been derived for A or B. We claim no bounded Hamiltonian, ghost freedom, stability, or causality. Geometric symmetries do not replace computing conserved charges. Higher order alone does not prove an Ostrogradsky theorem in a degenerate system.

\section{Standard limit and equivalence}\label{sec:p2-limits}
A recovers geodesic feasibility only exactly at zero; B has a nonuniform limit; C becomes nonrestrictive but has no dynamics. No convergence of solutions, initial data, or observables is shown. A/B remain in tension with pointwise equivalence principle.

\section{Observables}\label{sec:p2-observables}
No pointwise candidate yet maps data to $\kappa_0$. Excluding curves is not a quantitative prediction. C suggests RMS proper acceleration, but window, kernel, causality, and uncertainty model are absent.

\section{Decision and gate}\label{sec:p2-decision}
\texttt{NO\_GO\_NOT\_ESTABLISHED}: A and B are incomplete, not rejected. Paper III is scientifically deferred until at least one passes variation, constraints, stability, causality, standard-limit, and identifiability gates. C cannot serve as an implicit rescue.

\section{Limitations}\label{sec:p2-limitations}
We provide no complete equations, Hamiltonian analysis, worldline simulations, experimental bound, ALD, QED, gravity, or cosmology. Conclusions are conditional on fixed models.

\section{AI assistance}\label{sec:p2-ai}
AI assisted research, code, translation, and drafting; it is neither source nor coauthor. Metadata, algebra, interpretations, and conclusions require human responsibility.

\bibliographystyle{plain}
\bibliography{../../../references/library}
\end{document}
